\begin{table}[H]
\centering
{\tablefont
\begin{tabular}{@{}p{0.52\textwidth}l r@{}}
\toprule
Target family & Geographic level & Targets \\
\midrule
IRS Statistics of Income (\texttt{irs\_soi}) & National, state, district & 30{,}252 \\
Census population estimates (\texttt{census\_population}) & National, state & 936 \\
Medicaid and CHIP, CMS (\texttt{cms\_medicaid}) & National, state & 102 \\
ACA marketplace, CMS (\texttt{cms\_aca}) & State & 102 \\
SNAP, USDA (\texttt{usda\_snap}) & National, state & 52 \\
State income tax, Census STC (\texttt{state\_income\_tax}) & State & 44 \\
TANF, HHS ACF (\texttt{hhs\_acf\_tanf}) & National, state & 30 \\
Social Security OASDI, SSA (\texttt{ssa}) & National & 6 \\
JCT tax expenditures (\texttt{jct}) & National & 5 \\
CBO revenue projections (\texttt{cbo}) & National & 4 \\
Medicare, CMS (\texttt{cms\_medicare}) & National & 1 \\
\midrule
\textbf{Total} & & \textbf{31{,}534} \\
\bottomrule
\end{tabular}
}
\caption{Full materialized Populace target surface used in this study, by family and geographic level.
The surface contains 23{,}450 congressional-district targets, 7{,}615 state targets, and 469 national
targets.}
\label{tab:target_subset}
\tablenote{Counts are read from the full-surface run registry. In the headline experiment all targets,
including \texttt{cbo}, are fit and scored; holdout variants are supplemental diagnostics
(Section~\ref{sec:oos}). The source for each family is listed in Appendix~\ref{app:targets}.}
\end{table}
